\section{Protocolo experimental}

En esta sección se define el protocolo común para todos los experimentos, con el objetivo de garantizar comparabilidad entre enfoques y evitar \textit{data leakage}. El protocolo se aplica sobre el conjunto de modelado resultante de la Fase 1 \((3538, 16)\).

\subsection{Partición y validación}

Se adopta un esquema en dos niveles:

\begin{enumerate}
	\item \textbf{Partición externa estratificada}: separación inicial en \textit{train\_val} y \textit{test} mediante una única semilla fija (\texttt{random\_state = 42}). El conjunto de test queda bloqueado para evaluación final.
	\item \textbf{Validación interna estratificada}: dentro de \textit{train\_val}, uso de validación cruzada estratificada para seleccionar variables, hiperparámetros y umbral.
\end{enumerate}

Reglas operativas comunes:

\begin{itemize}
	\item Imputación, codificación y escalado se ajustan exclusivamente con datos de entrenamiento en cada pliegue.
	\item Cualquier selección de variables (RFE o importancia del árbol) se realiza dentro del marco de validación.
	\item El conjunto de test no interviene en la elección de transformaciones, hiperparámetros ni umbral.
\end{itemize}

\subsection{Validación de candidatas de feature engineering}

Las variables candidatas propuestas tras el EDA (tabla \texttt{fase2\_feature\_candidates.csv}) no se incorporan de forma automática. Su adopción se decide con una regla explícita de validación para evitar sobreajuste por diseño.

Para cada candidata (o bloque coherente de candidatas), se comparan dos configuraciones en validación cruzada sobre \textit{train\_val}:

\begin{itemize}
	\item \textbf{Modelo base}: sin la nueva transformación.
	\item \textbf{Modelo extendido}: incluyendo la transformación candidata.
\end{itemize}

El criterio de aceptación será:

\begin{enumerate}
	\item \textbf{Apartado orientado a Accuracy}: aceptar la candidata solo si mejora la Accuracy media de validación y no produce una caída relevante de Recall (caída absoluta mayor de 0.02).
	\item \textbf{Apartado orientado a Recall}: aceptar la candidata solo si mejora la Recall media de validación y no provoca un deterioro severo de Precision o PR-AUC (caída absoluta mayor de 0.02 en ambos indicadores).
	\item \textbf{Estabilidad}: en ambos apartados, la desviación estándar entre pliegues no debe aumentar de forma marcada respecto al modelo base (incremento superior al 20\% como señal de inestabilidad).
\end{enumerate}

Adicionalmente, si dos transformaciones ofrecen mejoras similares, se priorizará la alternativa más interpretable y con menor complejidad del pipeline.

\subsection{Selección de umbral}

La selección de umbral se realiza exclusivamente con predicciones de validación interna. Se evalúa una rejilla de umbrales y se guarda una tabla de métricas por umbral.

\begin{itemize}
	\item En el bloque de Accuracy, se selecciona el umbral que maximiza Accuracy manteniendo degradación de Recall controlada.
	\item En el bloque de Recall, se selecciona el umbral que prioriza captación de churn y se reporta explícitamente el coste en falsos positivos.
\end{itemize}

El umbral final elegido en validación se aplica una única vez sobre test.

\subsection{Métricas}

Las métricas comunes para todos los experimentos son:

\begin{itemize}
	\item Accuracy,
	\item Recall,
	\item Precision,
	\item F1-score,
	\item ROC-AUC,
	\item PR-AUC,
	\item matriz de confusión.
\end{itemize}

La métrica primaria cambia según el objetivo del apartado (Accuracy en el modelo 1 y Recall en el modelo 2), pero la interpretación final se realiza siempre de forma multivariable para evitar conclusiones sesgadas por un único indicador.
